\begin{center}
Problema F

A teia da aranha Tecelã

{\small Nome do arquivo: F.c, F.cpp, F.java, ou F.py}

{\large Timelimit: $1$}

{\tiny Por Douglas Cedrim}

\end{center}

	A Tecelã é uma aranha que vive antenada nas novas tecnologias que surgem pra ajudar no seu cotidiano. No seu dia a dia ela produz uma certa quantidade $X$ de fibra, que será utilizada como matéria prima nas teias que ela tece. Essa teia é muito importante pois permite, por exemplo, que ela consiga capturar insetos para sua própria alimentação.
	
%	Certo dia, ouvindo a conversa entre estudantes de Computação, ela ficou muito animada em conhecer o conceito de algoritmo. 
	%
	Ao visitar o prédio de Informática no IF Goiano, ela observou vários tubos onde poderia tecer sua teia para buscar alimento. Ela observou que todos eles tinham o formato cilíndrico, com as extremidades definidas por círculos.
	%	
	Decidiu então que só iria tecer uma teia se a quantidade $X$ de fibras que ela produziu, que é medida em $m^2$, fosse suficiente para cobrir toda a secção transversal da extremidade do tubo, e que precisaria bolar um algoritmo para isso. O problema é que os tubos têm tamanhos bem diferentes entre si.
	

	
	%Sendo assim, ela vai a um determinado tubo e percorre o contorno da secção transversal, produzindo uma teia para medir o seu comprimento. 
	Sendo assim, ela vai a um determinado tubo e percorre o contorno da secção transversal, produzindo uma teia que mede exatamente o comprimento $C$ da circunferência do cilindro. Considere que a espessura do cilindro é irrelevante para as contas. E também que o quanto ela usou de teia para a medição ($C$) não irá impactar na decisão.

	%Esse dado será utilizado como entrada no algoritmo para determinar o quanto de teia ela irá precisar para cobrir toda a extremidade da secção do tubo.
%	\vspace*{-2cm}
	\begin{figure}[h!]
		\centering
		\includegraphics[width=0.7\linewidth]{fig/f_aranha2.png}
		\caption{Esquema do algoritmo da aranha Tecelã}%, da espécie \textit{Parawixia bistriata}, muito comum no cerrado brasileiro.}
	
	\end{figure}

%	Considere que a medição da Tecelã é exata, e que a espessura do tubo é irrelevante para as contas.
	
	E aí, será que você pode ajudar a Tecelã a identificar se ela deve tecer a sua teia em um determinado tubo?

\vspace{0.3cm}
\textbf{Entrada}
Cada entrada contém duas medidas, separadas por um espaço: primeiro o montante $X$ de fibra produzido pela Tecelã, em $m^2$; e em seguida o comprimento $C$ da circunferência do tubo, medido com sua teia.

% A primeira linha contém um inteiro $X$ indicando o nível (iniciado em 0) da hierarquia onde estão os genes $(0 \leq X \leq 99)$. A segunda linha representa a quantidade $N$ de genes de um indivíduo $(1 \leq X \leq 100)$. A terceira linha é formada pela sequência de genes de um indivíduo, sendo que cada elemento está no intervalo de $0$ a $10000$.

\vspace{0.3cm}
\textbf{Saída}
O sistema deve retornar "SIM"\ caso a quantidade de fibra que ela possui seja suficiente para cobrir a secção da extremidade do tubo. Caso contrário, a resposta deverá ser "NAO". Não esqueça do terminador de quebra de linha.

%O sistema deve retornar a sequência de genes significativos seguido do caractere \#, todos separados por um único espaço, conforme o nível hierárquico $X$, indicado na primeira linha da entrada. Os filhos 1, 2, 3 e 4, devem ser representados da esquerda para direta. Se nenhum gene significativo for encontrado no nível indicado, a saída terá apenas o caractere \#.

\begin{table}[ht]
    \centering
   % \small
    \begin{tabular}{|p{12cm}|p{3cm}|}
        \hline
        \begin{center}
            \vspace{-0.3cm}
            \textbf{Exemplos de Entrada}
            \vspace{-0.3cm}
        \end{center} 
         &  
        \begin{center}
            \vspace{-0.3cm}
            \textbf{Exemplos de Saída}
            \vspace{-0.3cm}
        \end{center} \\ \hline
		    \texttt{3.4} \texttt{6.3}  & \texttt{SIM} \\
         \hline
   		    \texttt{3.0} \texttt{6.3}  & \texttt{NAO} \\
         \hline    
   		    \texttt{0.09} \texttt{1}  & \texttt{SIM} \\
%         \hline    
         %\texttt{2} & \texttt{2501 25 38 78 \#} \\  
         %\texttt{10} & \texttt{} \\  
         %\texttt{5 98 78 50 7 38 25 61 1579 2501} & \texttt{} \\
         \hline    
    \end{tabular}
    \caption{Exemplos de entradas e saídas}

\end{table}

%\normalsize

\newpage

\begin{center}
\noindent
\lstinputlisting{abobora_24/F/F4.cpp}
\end{center}